\section*{第 \ref{ch:b2:orbital-forcing} 章——轨道强迫与冰期}

\begin{solution}{exo:b2:orbital-forcing:1}
(a) $1.0167/0.9833 = 1.0340$：远日点比近日点远 \qty{3.4}{\percent}。
(b) $a(1 + e) - a(1 - e) = 2ea = 0.0334\,a$，为 $a$ 的 \qty{3.34}{\percent}，与
通常引用的 3.4\% 相近。(c) 日光强度与距离的平方成反比：$1.0340^2 = 1.069$，
近日点处的日光约多 \qty{7}{\percent}。
\end{solution}

\begin{solution}{exo:b2:orbital-forcing:2}
(a) 现今 $0.017236 \times \sin\ang{101.37} = 0.017236 \times 0.9804 = 0.01690$；\num{10000} 年前
$0.019971 \times \sin\ang{293.86} = 0.019971 \times (-0.9145) = -0.01826$；\num{6000} 年前
$0.019249 \times \sin\ang{359.99} \approx 0$，与表中数值一致。(b) 太阳在近日点的黄经为
$\varpi + \ang{180}$：现今为 \ang{281.37}，
\num{10000} 年前为 \ang{473.86} $-$ \ang{360} $=$ \ang{113.86}，\num{6000} 年前为 \ang{179.99}。(c) \num{10000} 年前，近日点
落在六月至点之后 \ang{24}，是三者中最接近六月至点的一次，此时北半球为夏季；
\num{6000} 年前近日点落在九月分点，此时\omterm{def:b2:orbital-forcing:elements}{气候岁差}为零。
\end{solution}

\begin{solution}{exo:b2:orbital-forcing:3}
(a) 频率相加：$1/25.77 + 1/112 = 0.03880 + 0.00893 = \qty{0.04773}{ka^{-1}}$，对应周期为
$1/0.04773 = 21.0$ 千年。(b) 结果落在主要周期的范围之内，
即 19.0 至 23.7 千年。单一数值无法重现这些周期：轨道的转动本身
是若干频率不同的运动之和，每一项都与地轴的岁差结合，
给出一个\omterm{def:b2:orbital-forcing:elements}{气候岁差}周期；平均值即
通常引用的 \num{23000} 年。
\end{solution}

\begin{solution}{exo:b2:orbital-forcing:4}
$\sqrt{1 - 0.0167^2} = 0.99986$。对于 \ang{22.1}：$1360 \times 0.37622/(\pi \times 0.99986) =
\qty{162.9}{W.m^{-2}}$；对于 \ang{24.5}：$1360 \times 0.41469/(\pi \times 0.99986) =
\qty{179.5}{W.m^{-2}}$。比值为 $\sin\ang{24.5}/\sin\ang{22.1} = 1.102$：极地的年平均值
在两个极端之间增加 \qty{10.2}{\percent}，约为 \qty{16.7}{W.m^{-2}}。
\end{solution}

\begin{solution}{exo:b2:orbital-forcing:5}
(a) 在六月至点，$(a/r)^2 = \bigl((1 - e\sin\varpi)/(1 - e^2)\bigr)^2$。现今 $e\sin\varpi
= 0.016898$，$(0.983102/0.999703)^2 = 0.9671$，对比 $1 - 0.033796 = 0.9662$；
\num{116000} 年前 $e\sin\varpi = 0.043939$，$(0.956061/0.998065)^2 = 0.9176$，对比 $1 -
0.087878 = 0.9121$。(b) 在十二月至点，$(a/r)^2 = \bigl((1 + e\sin\varpi)/(1 -
e^2)\bigr)^2$：现今为 $1.0347$，对比 $1.0338$；\num{116000} 年前为 $1.0940$，对比 $1.0879$。
(c) 略去的项为 $e^2$ 量级：$(1 - x)^2/(1 - e^2)^2 \approx 1 - 2x + x^2 + 2e^2$，
其中 $x = e\sin\varpi$，此处约为 $3e^2$。现今其值为 $3 \times 0.0172^2 = 0.0009$，\num{116000} 年前为 $3
\times 0.0440^2 = 0.0058$，与所得差值 $0.0009$ 和
$0.0055$ 相近：偏心率增大至约两倍半时，近似的误差
约为原来的六倍。
\end{solution}

\begin{solution}{exo:b2:orbital-forcing:6}
在六月至点，$\delta = \varepsilon = \ang{23.446}$。$\cos h_0 = -\tan\ang{65}\tan\ang{23.446}
= -2.14451 \times 0.43369 = -0.93006$，故 $h_0 = \ang{158.44} = \qty{2.7654}{rad}$。日照时长为
$2 \times 158.44/15 = 21.13$~小时（地球每小时转过 \ang{15}）。取 $\sin\ang{65} = 0.90631$，
$\cos\ang{65} = 0.42262$，$\sin\delta = 0.39788$，$\cos\delta = 0.91744$，及 $\sin h_0 = 0.36742$，
括号中的值为 $2.7654 \times 0.90631 \times 0.39788 + 0.42262 \times 0.91744 \times 0.36742 = 0.99720
+ 0.14246 = 1.13966$。在平均距离处 $\bar Q = (1360/\pi) \times 1.13966 = 432.90 \times 1.13966
= \qty{493.4}{W.m^{-2}}$；乘以\cref{exo:b2:orbital-forcing:5}中的距离因子，$493.4 \times
0.9671 = \qty{477.1}{W.m^{-2}}$。作为校核，对北纬 \ang{60} 作同样计算，得 \qty{475.96}{W.m^{-2}}，表中数值为 \num{475.95}。
\end{solution}

\begin{solution}{exo:b2:orbital-forcing:7}
(a) 三月分点为 $\lambda = 0$，九月分点为 $\lambda = \pi$。取 $\sqrt{1 - e^2}
\approx 1$，时长为 $\frac{T}{2\pi}\int_0^{\pi}\bigl(1 + 2e\cos(\lambda -
\varpi)\bigr)\,\mathrm d\lambda = \frac{T}{2\pi}\bigl(\pi + 2e[\sin(\lambda - \varpi)]_0^{\pi}\bigr)
= \frac{T}{2\pi}(\pi + 4e\sin\varpi) = \frac T2 + \frac{2T}{\pi}e\sin\varpi$。(b) 另一半年较短，相差同样的量，
故两者之差为 $(4T/\pi)\,e\sin\varpi = (4 \times
365.2422/\pi) \times 0.016898 = 465.04 \times 0.016898 = 7.86$~天。(c) 春季与夏季合计
比秋季与冬季多 $4.5 + 3 = 7.5$~天；所引两个差值均为取整数值，
二者的一致性在此取整误差范围之内。
\end{solution}

\begin{solution}{exo:b2:orbital-forcing:8}
(a) 对 $e = 0.0046$：$\sqrt{1 - e^2} = 0.999989$，$340/0.999989 = \qty{340.004}{W.m^{-2}}$。对
$e = 0.0565$：$\sqrt{1 - 0.003192} = 0.998403$，$340/0.998403 = \qty{340.544}{W.m^{-2}}$。(b) 两者之差
为 \qty{0.54}{W.m^{-2}}，即 \qty{0.16}{\percent}。(c) 北纬 \ang{65} 处七月中旬的日射量
变化了 $497 - 388 = \qty{109}{W.m^{-2}}$，约占其平均值的 \qty{25}{\percent}，约为
全球年平均值变化的两百倍：轨道对气候的作用在于使日光在季节与纬度之间重新分配，
而不是改变其总量。
\end{solution}

\begin{solution}{exo:b2:orbital-forcing:9}
取 $\omega = 2\pi/P$，$\tau = \qty{10}{ka}$：对于 $P = 23$，$\omega\tau = 2.732$，相对振幅
$1/\sqrt{1 + 7.464} = 0.344$，滞后 $\arctan(2.732)/\omega = 1.2199/0.2732 = 4.47$~千年；
对于 $P = 41$，$\omega\tau = 1.532$，振幅 $0.546$，滞后 $0.9926/0.1532 = 6.48$~千年；对于 $P =
100$，$\omega\tau = 0.628$，振幅 $0.847$，滞后 $0.5610/0.0628 = 8.93$~千年。冰盖偏向于慢速强迫：
它跟随 \num{100000} 年周期的 \qty{85}{\percent}，却只跟随岁差周期的
三分之一，因为在快速强迫反向之前，它没有足够的时间增长或消退。
它起着低通滤波器的作用。
\end{solution}

\begin{solution}{exo:b2:orbital-forcing:10}
(a) 该方程是线性的：若 $V_1$ 对应 $A_1\sin\omega_1t$，$V_2$ 对应
$A_2\sin\omega_2t$，则 $V_1 + V_2$ 对应二者之和。由\cref{prop:b2:orbital-forcing:lag}可知，$V_1$ 与
$V_2$ 分别为频率是 $\omega_1$ 与 $\omega_2$ 的正弦函数（瞬变过程消失之后），
因此响应中不含其他频率。(b) 取 $V \approx B_1\sin(\omega_1t - \phi_1) +
B_2\sin(\omega_2t - \phi_2)$,
\[
\begin{aligned}
  V^2 &= \tfrac12(B_1^2 + B_2^2) - \tfrac12B_1^2\cos(2\omega_1t - 2\phi_1) -
  \tfrac12B_2^2\cos(2\omega_2t - 2\phi_2)\\
  &\quad + B_1B_2\bigl[\cos\bigl((\omega_1 - \omega_2)t - \phi_1 + \phi_2\bigr) -
  \cos\bigl((\omega_1 + \omega_2)t - \phi_1 - \phi_2\bigr)\bigr],
\end{aligned}
\]
因此项 $kV^2$ 使方程在差频与和频处受迫，以及
在倍频处受迫。(c) $1/(1/19.0 - 1/23.7) = 1/(0.05263 - 0.04219) = 95.8$ 千年；由 19.0
与 22.4 组成的一对给出 125 千年。岁差项的差频落在
偏心率频带内，这是必然的，因为\omterm{def:b2:orbital-forcing:elements}{气候岁差}即为 $e\sin\varpi$。非线性的
冰盖因此可以从岁差中析出一个 \num{100000} 年的节律；线性的则不能。
\num{100000} 年周期是非线性行为的证据，不过
仅凭本计算无法确定究竟是哪一种。
\end{solution}

\begin{solution}{exo:b2:orbital-forcing:11}
(a) 与\cref{prop:b2:orbital-forcing:annual}的证明一样，$(a/r)^2\,\mathrm dt =
T\,\mathrm d\lambda/(2\pi\sqrt{1 - e^2})$，且 $\sin\delta = \sin\varepsilon\sin\lambda$ 在
$0 < \lambda < \pi$ 时为正。于是 $E = \frac{S_0T\sin\varepsilon}{2\pi\sqrt{1 - e^2}}\int_0^{\pi}\sin\lambda
\,\mathrm d\lambda = S_0T\sin\varepsilon/(\pi\sqrt{1 - e^2})$。(b) $T = 365.25 \times 86400 =
\qty{3.1558e7}{s}$，$S_0T/\pi = \qty{1.3661e10}{J.m^{-2}}$；当 \ang{22.1} 时，$E = 1.3661 \times 10^{10}
\times 0.37622/0.99986 = \qty{5.14}{GJ.m^{-2}}$；当 \ang{24.5} 时，$\qty{5.67}{GJ.m^{-2}}$，
多 \qty{10}{\percent}。(c) 岁差决定同样的能量在时间上如何分布：近日点的夏季
较短而更强烈，远日点的夏季较长而较弱。融化
只在高于 \qty{0}{\celsius} 时开始，因此短促而强烈的夏季可以融化
比接收同样总量的漫长而凉爽的夏季融化得更多。
\end{solution}

\begin{solution}{exo:b2:orbital-forcing:12}
(a) $3.93\log_2(279.4/189.7) = 3.93\log_2 1.4729 = 3.93 \times 0.5586 = \qty{2.20}{W.m^{-2}}$。(b) 最小的
增幅是从 194 升至 \qty{266}{ppm}：$3.93\log_2 1.3711 = 3.93 \times 0.4554 =
\qty{1.79}{W.m^{-2}}$；最大增幅则从 188 升至 \qty{282}{ppm}：$3.93\log_2 1.5 = 3.93 \times 0.5850 =
\qty{2.30}{W.m^{-2}}$。(c) 二氧化碳的强迫相当于
全球年日射量最大变化（\qty{0.54}{W.m^{-2}}）的四倍，且作用于全球各处、全年各时；
而北纬 \ang{65} 的夏季日射量约变化 \qty{109}{W.m^{-2}}，只作用于一个纬度、一个季节。轨道是
触发因素，通过北方冰盖的夏季起作用；二氧化碳与
\omterm{def:b2:precambrian-earth:ice-albedo}{冰–反照率反馈}一起，是把这一变化扩散到全球的放大器。
\end{solution}

\begin{solution}{pb:b2:orbital-forcing:1}
\textbf{1.} 则 $0.0440 \times \sin\ang{92.7} = 0.0440 \times 0.99889 = 0.04395$；现在 $0.0172 \times
\sin\ang{101.4} = 0.0172 \times 0.98027 = 0.01686$，约小 2.6 倍。
\textbf{2.} 则 $1 - 0.0440 = 0.956$ 与 $1 + 0.0440 = 1.044$；现在为 $0.9828$ 与 $1.0172$。
\textbf{3.} 当 $\lambda = \varpi + \ang{180}$ 时：当时为 \ang{272.7}，现在为 \ang{281.4}，两者都接近
冬至（$\lambda = \ang{270}$）附近；当时地球在该日后不到三天
即经过近日点，远日点在夏至。
\textbf{4.} 则 $(1.044/0.956)^2 = 1.0921^2 = 1.193$，现在 $(1.0172/0.9828)^2 = 1.0350^2 = 1.071$：
近日点的日射量当时多 \qty{19}{\percent}，现在多 \qty{7}{\percent}。
\textbf{5.} 更小，为 \ang{22.52}，而现在为 \ang{23.45}：高纬度地区的夏季太阳位置更低，
高纬度地区在夏季及全年得到的日射量都更少。
\textbf{6.} 则 $340/\sqrt{1 - 0.001936} = 340/0.999032 = \qty{340.33}{W.m^{-2}}$；现在 $340/0.999852
= \qty{340.05}{W.m^{-2}}$。两者之差 \qty{0.28}{W.m^{-2}}，不足千分之一，
可以忽略。
\textbf{7.} 当 $\lambda = \ang{90}$ 时，$\cos(\ang{90} - \varpi) = \sin\varpi$，
见\cref{prop:b2:orbital-forcing:distance}。
\textbf{8.} 则 $(1 - 0.001936)/(1 - 0.04395) = 0.998064/0.956049 = 1.04395$；现在 $(1 -
0.000296)/(1 - 0.01686) = 0.999704/0.983139 = 1.01685$.
\textbf{9.} 则 $1/1.04395^2 = 0.91758$；现在 $1/1.01685^2 = 0.96714$。
\textbf{10.} 则 $1 - 2 \times 0.04395 = 0.91210$；现在 $1 - 2 \times 0.01686 = 0.96628$：近似
的误差当时为 $0.0055$，现在为 $0.0009$。
\textbf{11.} 则 $1 - 0.91758 = \qty{8.2}{\percent}$ 偏弱，现在则为 \qty{3.3}{\percent}。
\textbf{12.} 至日时 $\delta = \varepsilon$：当时为 \ang{22.52}，现在为 \ang{23.45}。
\textbf{13.} 则 $\cos h_0 = -2.14451 \times 0.41462 = -0.88916$，$h_0 = \ang{152.77} =
\qty{2.6663}{rad}$；现在 $\cos h_0 = -2.14451 \times 0.43378 = -0.93023$，$h_0 = \ang{158.47} =
\qty{2.7659}{rad}$。
\textbf{14.} 每小时 $2h_0/\ang{15}$：当时为 $20.37$~小时，现在为 $21.13$~小时。
\textbf{15.} 则 $2.6663 \times 0.90631 \times 0.38301 + 0.42262 \times 0.92375 \times 0.45759 =
0.92554 + 0.17864 = 1.10418$（未舍入时为 \num{1.10417}），且 $432.90 \times 1.10417 =
\qty{478.0}{W.m^{-2}}$；现在 $2.7659 \times 0.90631 \times 0.39795 + 0.42262 \times 0.91741 \times
0.36697 = 0.99757 + 0.14228 = 1.13985$（未舍入时为 \num{1.13982}），且 $432.90 \times 1.13982 = \qty{493.4}{W.m^{-2}}$。
\textbf{16.} 则 $478.0 \times 0.91758 = \qty{438.6}{W.m^{-2}}$；现在 $493.4 \times 0.96714 =
\qty{477.2}{W.m^{-2}}$，比当时多 \qty{38.6}{W.m^{-2}}。
\textbf{17.} 七月中旬太阳的黄经为 \ang{120}：$\sin\delta = \sin\varepsilon\sin\ang{120}$，
赤纬约为 \ang{20}，而非 \ang{23}，因此白昼更短、太阳位置更低。在
北纬 \ang{65} 处，日平均日射量在夏至附近最大，此时赤纬最大。
\textbf{18.} 按今天的距离计，仅黄赤交角变化的影响为 $(493.4 - 478.0) \times
0.96714 = \qty{14.9}{W.m^{-2}}$；若用旧的黄赤交角，仅距离的影响为 $478.0 \times
(0.96714 - 0.91758) = \qty{23.7}{W.m^{-2}}$。两者相加为 \qty{38.6}{W.m^{-2}}；距离的
作用约占五分之三。
\textbf{19.} 骤然变化之后，冰盖已达到该变化的 $1 - \mathrm e^{-t/\tau}$：
$1 - \mathrm e^{-0.5} = 0.39$ 经过 \qty{5}{ka} 时，$1 - \mathrm e^{-1} = 0.63$ 经过 \qty{10}{ka} 时，
$1 - \mathrm e^{-2} = 0.86$ 经过 \qty{20}{ka} 时。
\textbf{20.} $\omega\tau = 2\pi \times 10/23 = 2.732$：冰盖对强迫的跟随程度为 $1/\sqrt{1 + 2.732^2} =
0.34$，滞后 $\arctan(2.732)/\omega = 4.5$~ka。
\textbf{21.} 半个岁差周期，约 11 至 \qty{12}{ka}，约等于一个响应时间：在
冰盖推进到三分之二的路程之前，日射量便已转向，把它拉了回去。
\textbf{22.} 在线性模型中，冰盖增长与消融的速率相同，
只要强迫大小相同：上升与下降的强迫给出镜像式的响应。先缓慢积累、后
快速崩塌，需要一种取决于状态或变化方向的响应，即
非线性的响应。
\textbf{23.} 二氧化碳从\omterm{def:b2:orbital-forcing:glacials}{冰期}到\omterm{def:b2:orbital-forcing:glacials}{间冰期}使强迫增加了 $3.93\log_2(266/194) = 1.8$ 至 $3.93\log_2(282/188) =
\qty{2.3}{W.m^{-2}}$，是覆盖全球和
全年的强迫。二氧化碳在大西洋环流减弱时被释放，它把北方夏季的变化
带到整个星球，并与缩小中的冰盖反照率一起将其放大。
\textbf{24.} 末次\omterm{def:b2:orbital-forcing:glacials}{间冰期}结束时，夏至日平均日射量在
北纬 \ang{65} 处约为 \textbf{\qty{439}{W.m^{-2}}}，比今天少 \qty{39}{W.m^{-2}}，且为
该周期中的最低值。
\end{solution}
